\documentclass[a4paper]{article}

\usepackage[fontset=fandol]{ctex}
\usepackage{amsmath, amssymb}
\usepackage{graphicx}
\usepackage[margin=2.5cm]{geometry}
\usepackage[most]{tcolorbox}
\usepackage{etoolbox}
\usepackage{listings}
\usepackage{hyperref}
\usepackage{booktabs}
\usepackage{subcaption}
\usepackage{float}
\usepackage{tikz}
\IfFileExists{pgfplots.sty}{
    \usepackage{pgfplots}
    \pgfplotsset{compat=1.18}
}{}

\newtcolorbox{knowledgebox}[1]{
    enhanced, colback=blue!5!white, colframe=blue!75!black, colbacktitle=blue!75!black,
    coltitle=white, fonttitle=\bfseries, title=#1, attach boxed title to top left={yshift=-2mm, xshift=2mm},
    boxrule=1pt, sharp corners
}

\newtcolorbox{importantbox}[1]{
    enhanced, colback=yellow!10!white, colframe=yellow!80!black, colbacktitle=yellow!80!black,
    coltitle=black, fonttitle=\bfseries, title=#1, sharp corners
}

\newtcolorbox{warningbox}[1]{
    enhanced, colback=red!5!white, colframe=red!75!black, colbacktitle=red!75!black,
    coltitle=white, fonttitle=\bfseries, title=#1, sharp corners
}

\lstset{
    language=Python,
    basicstyle=\ttfamily\small,
    keywordstyle=\color{blue},
    stringstyle=\color{red!60!black},
    commentstyle=\color{green!60!black},
    breaklines=true,
    frame=single,
    numbers=left,
    numberstyle=\tiny\color{gray},
    captionpos=b,
    extendedchars=false
}

\newcommand{\notetitle}{CS231n Lecture 8: Attention and Transformers}
\newcommand{\noteauthors}{Codex Harness}
\newcommand{\notedate}{2026-04-15}
\newcommand{\videochannel}{Stanford Online / CS231n}
\newcommand{\videopublishdate}{2025-04-24}
\newcommand{\videoduration}{01:07:16}
\newcommand{\videourl}{https://www.youtube.com/watch?v=RQowiOF_FvQ}
\newcommand{\videocoverpath}{cover.jpg}

\begin{document}

\begin{titlepage}
\centering
{\Large 课程笔记\par}
\vspace{1.2cm}
{\huge\bfseries \notetitle\par}
\vspace{0.8cm}
{\large \noteauthors\par}
\vspace{0.3cm}
{\large \notedate\par}
\vspace{1.2cm}
\includegraphics[width=0.82\textwidth,height=0.45\textheight,keepaspectratio]{\videocoverpath}\par
\vfill
\begin{tcolorbox}[width=0.9\textwidth, colback=black!2!white, colframe=black!60, sharp corners]
\textbf{视频作者/频道}：\videochannel\par
\textbf{发布日期}：\videopublishdate\par
\textbf{视频时长}：\videoduration\par
\textbf{视频链接}：\href{\videourl}{\nolinkurl{\videourl}}
\end{tcolorbox}
\end{titlepage}

\tableofcontents
\newpage

\section{从序列到序列到注意力: 为什么瓶颈会出现}
% source_units: cu_08_02
上一讲的 RNN 告诉我们，序列建模可以通过 hidden state 把历史信息递推下去。但这讲一开始就把一个关键问题摆到了台面上：如果输入和输出都很长，只靠一个固定长度的上下文向量，decoder 往往会被迫在“记住一切”和“立刻生成下一步”之间做不可能的折中。

在 sequence-to-sequence（seq2seq）模型中，经典 encoder-decoder 结构会先把整条输入序列编码成最后一个状态，再把这个状态交给 decoder 逐步生成输出。这个方案的短板很明显：
\begin{itemize}
\item 输入越长，最后一个状态承载的信息越拥挤。
\item decoder 在第 $t$ 步生成时，无法有选择地回看输入中的局部片段。
\item 某些局部对齐关系，例如翻译中的词对齐或视觉描述中的局部区域对齐，很难通过单个向量精确表达。
\end{itemize}

\begin{figure}[H]
\centering
\includegraphics[width=0.86\textwidth]{pdf_pages/page-13.png}
\caption{Attention 最初是为了解决 encoder-decoder 的固定瓶颈：decoder 需要按步访问 encoder 中不同位置的信息。}
\end{figure}

因此，课程把 attention 定义为一种新原语：与其把所有信息预先压缩成一个向量，不如让当前步骤主动去“检索”它需要的上下文。

\subsection{本章小结}
attention 的出发点不是替换一层算子，而是推翻“所有历史必须压进单一上下文向量”的假设。它让模型从固定压缩转向按需读取。

\section{Cross-attention: 用可学习检索替代固定长度压缩}
% source_units: cu_08_03, cu_08_04
如果把 decoder 当前状态记作 query，把 encoder 的隐藏状态记作 keys 和 values，那么 attention 的基本步骤可以写成三步：

\[
e_{i,j} = \frac{q_i^\top k_j}{\sqrt{d_k}}
\]
\[
a_{i,j} = \operatorname{softmax}_j(e_{i,j})
\]
\[
y_i = \sum_j a_{i,j} v_j
\]

其中：
\begin{itemize}
\item $q_i$：第 $i$ 个输出位置对应的 query。
\item $k_j$：第 $j$ 个输入位置对应的 key。
\item $v_j$：第 $j$ 个输入位置对应的 value。
\item $e_{i,j}$：query 与 key 的匹配分数。
\item $a_{i,j}$：归一化后的注意力权重。
\item $y_i$：当前输出位置聚合得到的上下文向量。
\end{itemize}

这个形式揭示了 attention 的本质：\textbf{当前步骤不再被动接受一个固定上下文，而是显式地在所有输入位置上分配概率质量。}

当有多个 query 一起计算时，上面的公式可以矩阵化：
\[
A = \operatorname{softmax}\!\left(\frac{QK^\top}{\sqrt{d_k}}\right), \qquad
Y = AV
\]

这里 $A$ 的每一行或每一列（取决于实现约定）都表示一个 query 对所有 key 的分布。课程在 slides 里特意把这一层层展开，是为了让读者看到：attention 并不神秘，它就是“先算相似度，再做 softmax，再用它对 value 做加权和”。

\begin{figure}[H]
\centering
\includegraphics[width=0.86\textwidth]{pdf_pages/page-48.png}
\caption{Attention matrix 可以理解为 query 对 key 的检索权重表，而输出是这些权重对 values 的线性组合。}
\end{figure}

\begin{importantbox}{为什么要除以 $\sqrt{d_k}$}
课程强调，如果点积的尺度太大，softmax 容易进入过饱和区间，梯度就会变差。缩放因子 $\sqrt{d_k}$ 的角色，就是把分数维持在更稳定的数值范围内。
\end{importantbox}

\subsection{本章小结}
cross-attention 允许 decoder 针对不同输出位置读取不同输入位置，从而把“上下文”从固定向量改造成动态检索结果。

\section{Self-attention 的矩阵视角}
% source_units: cu_08_05, cu_08_06
当 query、key、value 都来自同一条输入序列时，attention 就变成了 \textbf{self-attention}。设输入为 $X \in \mathbb{R}^{N \times D}$，那么：

\[
Q = XW_Q,\qquad K = XW_K,\qquad V = XW_V
\]
\[
A = \operatorname{softmax}\!\left(\frac{QK^\top}{\sqrt{d_k}}\right),\qquad
Y = AV
\]

这套写法非常关键，因为它直接把序列内部的所有 pairwise interaction 统一成矩阵乘法。课程里反复展示 attention weights，目的就是让学生从“递归式记忆”切换到“全局关系图”的视角：每个 token 都可以直接和其它 token 建立联系。

\begin{knowledgebox}{与 RNN 的根本区别}
RNN 依赖状态链逐步传递信息；self-attention 则让任意两个位置在单层里直接交互。前者是路径长度随距离增长，后者是路径长度近似常数，但代价是 $N \times N$ 的关系矩阵。
\end{knowledgebox}

这也解释了为什么课程同时强调 self-attention 的能力与代价：
\begin{itemize}
\item 好处：长程依赖不需要穿过很长的状态链。
\item 代价：时间和显存都随序列长度平方增长。
\end{itemize}

\subsection{本章小结}
self-attention 把“序列处理”变成了“关系矩阵计算”。这让模型更容易捕捉远程依赖，但也引入了二次复杂度的现实约束。

\section{位置编码、mask 与多头机制}
% source_units: cu_08_07, cu_08_08, cu_08_09
单纯的 self-attention 有一个严重问题：如果不加入位置信息，模型看到的只是一组无序 token，交换 token 顺序并不会改变 pairwise similarity 的集合。因此，课程把 \textbf{位置编码（positional encoding）} 作为 Transformer 的必要补丁，而不是锦上添花。

另一项关键机制是 \textbf{masked self-attention}。在自回归生成里，第 $t$ 个位置只能访问自己之前的 token，不能偷看未来。因此 mask 的角色，是在 softmax 之前把不该看的位置强行压成无穷小概率。

\begin{figure}[H]
\centering
\includegraphics[width=0.86\textwidth]{pdf_pages/page-75.png}
\caption{多头注意力和 mask 一起定义了 Transformer 的上下文可见性与表征分工。}
\end{figure}

多头注意力（multi-head self-attention）则进一步回答了另一个问题：如果一套 Q/K/V 投影只能在一个子空间里建模关系，是否会太单一？课程给出的答案是并行运行 $H$ 份注意力，再把结果拼接起来：
\[
\operatorname{MHA}(X) = \operatorname{Concat}(\text{head}_1,\dots,\text{head}_H)W_O
\]

每个 head 都有自己的一组投影矩阵，因此它们可以分别捕捉不同粒度、不同语义的依赖。工程上看，多头机制把大一次投影拆成多个小子空间；建模上看，它避免了“所有关系都挤在一张 attention map 上”。

\subsection{本章小结}
位置编码解决顺序感缺失，mask 解决自回归信息泄露，多头机制解决单一注意力空间过于拥挤的问题。三者合在一起，self-attention 才真正可用。

\section{Transformer block、训练目标与推理方式}
% source_units: cu_08_10, cu_08_11
Transformer 真正强大的地方，不只是 self-attention 这一个层，而是把注意力、前馈网络、残差连接和归一化组织成了可重复堆叠的 block。

\begin{lstlisting}[caption={简化版 Transformer block / self-attention 伪代码}]
def self_attention_block(X):
    Q = X @ W_Q
    K = X @ W_K
    V = X @ W_V
    A = softmax((Q @ K.T) / sqrt(d_k), dim=-1)
    Y = A @ V
    X = layer_norm(X + Y)
    Z = mlp(X)
    return layer_norm(X + Z)
\end{lstlisting}

课程里专门用“Self-Attention is Four Matrix Multiplies!” 来强调：从实现角度看，Transformer 的主体可以高度规约成矩阵乘法和简单逐元素操作，这正是它后来能在 GPU/TPU 上被大规模并行化的原因之一。

\begin{figure}[H]
\centering
\includegraphics[width=0.9\textwidth]{pdf_pages/page-90.png}
\caption{Transformer block 把注意力、残差、归一化与 MLP 拼成标准模块，使模型可以稳定地向深层堆叠。}
\end{figure}

训练目标方面，课程延续的是 next-token prediction：给定前缀，预测下一个 token 的分布。这一目标可以落在 encoder-decoder、decoder-only 等多种结构上；区别主要在于：
\begin{itemize}
\item encoder-decoder 更适合“输入序列 $\rightarrow$ 输出序列”的映射。
\item decoder-only 更适合统一的自回归建模。
\item masked self-attention 保证训练时目标与推理时可见信息一致。
\end{itemize}

\subsection{本章小结}
Transformer 的成功来自“新原语 + 可扩展 block 结构 + 适合并行硬件”的组合，而不是单一公式。attention 提供表达力，block 提供工程可扩展性。

\section{Vision Transformer 与近期结构 tweak}
% source_units: cu_08_12
本讲最后把注意力正式带回视觉任务。Vision Transformer（ViT）的关键动作是把图像切成 patch，并把每个 patch 当成一个 token 输入 Transformer。这样一来，NLP 里的 token sequence 被替换成了视觉 patch sequence，但 block 的核心形式几乎不变。

\begin{figure}[H]
\centering
\includegraphics[width=0.86\textwidth]{pdf_pages/page-105.png}
\caption{ViT 把图像 patchify 成 token，再用 Transformer 建模全局视觉关系。}
\end{figure}

与此同时，课程还快速回顾了若干“Transformer 体系结构没有根本变化，但细部实现持续演进”的 tweak：
\begin{itemize}
\item \textbf{Pre-Norm}：把归一化前置，提高深层训练稳定性。
\item \textbf{RMSNorm}：以更轻量的方式做尺度控制。
\item \textbf{SwiGLU}：在 MLP 中使用更强的门控非线性。
\item \textbf{MoE}：让部分前馈子模块稀疏激活，提高参数规模与算力利用的折中。
\end{itemize}

这些 tweak 说明，Transformer 的中心骨架很稳定，但在数值稳定性、计算效率与容量扩展上，仍然有很多工程优化空间。

\subsection{本章小结}
ViT 说明 attention 并不局限于文本，只要能把输入离散成 token 序列，就可以把 Transformer 迁移过去。而近期 tweak 则说明，成熟架构的演进更多体现在稳定性和效率。

\section{总结与延伸}
这讲完成了一个非常重要的转折：从“用状态链传递历史”的 RNN 视角，切换到“用可学习检索和关系矩阵处理上下文”的 attention / Transformer 视角。核心收获可以概括为三点：

\begin{itemize}
\item attention 解决了 encoder-decoder 的固定瓶颈，让模型能按需回看输入。
\item self-attention 把序列内部依赖统一成矩阵运算，为并行训练打下了基础。
\item Transformer block 把注意力、MLP、残差和归一化组织成了可堆叠、可扩展、可迁移到视觉的标准模块。
\end{itemize}

对计算机视觉而言，ViT 的意义不只是“把 Transformer 搬到图像上”，而是提示我们：一旦 tokenization 设计得当，很多视觉任务也可以享受 attention 带来的全局建模能力。下一讲进入 detection、segmentation 和 model interpretation 时，这种“全局关系建模”的思想会进一步和空间预测任务结合起来。

\end{document}
